\section{Marked determinant continuations and an adaptive exact fallback}
\label{sec:determinant-continuations}

Report 26 supplies a stronger observation of a partially scanned complex:
complete each genuine relative direct summand with the actual remaining
diagram, reduce at a marked point, and retain its Euler characteristic in
each quantum degree modulo four.  Its norm can detect knottedness before
the whole scan closes.  The maintained implementation is
\code{fastunknot/shadow\_scan.py}, selected by \code{backend="shadow"}.
It evaluates actual completed diagrams by exact integer determinants.
The report's reusable terminal kernel is verified separately; a new
geometric producer for that kernel is still a research prototype.
Neither observation gives a general bound on when a useful summand appears.

\subsection{Why the marked completion is necessary}

At a proper prefix of a crossing order, choose a point on an edge incident
to the last unprocessed crossing.  The remaining tangle retains this point
through every prefix observation.  The marked circle carries the action
of $A=\F[x]/(x^2)$.  The relative cobordism maps, cancellations, and their
explicit chain homotopies commute with this action.  Completing the relative
complex and applying the marked quotient therefore preserves these chain
homotopies: apply the quotient to each identity $fg-1=dh+hd$.
No assertion that the quotient functor is exact or flat is needed.

Suppose the maintained relative complex splits as whole differential
components $C_a$, with multiplicities $m_a$.  After the actual suffix is
attached and the marked quotient is applied, it follows that
\[
 \Kh(K;\F)\cong\bigoplus_a
 H\bigl(\widetilde{\operatorname{cl}}(C_a)\bigr)^{\oplus m_a},
\]
up to the recorded grading shifts.  This is a statement about components
of a genuine relative knot complex.  An arbitrary graded matrix satisfying
$d^2=0$ is insufficient evidence of this provenance.

There is also a crucial endpoint restriction.  After the final crossing
is processed, individual unmarked scalar generators need not be free
marked $A$-summands.  Dividing the Euler value of each such generator by
$q+q^{-1}$ would be invalid.  The observer rejects the final stage, and
the complete fallback instead uses the established unreduced rank test.

\subsection{A four-coordinate lower bound}

For a finite bigraded completed complex $X$, write
\[
 \sigma_r(X)=\sum_{h,\;j\equiv r\ (4)}(-1)^h\dim_{\F} X^{h,j},
 \qquad r=0,1,2,3.
\]
This vector is unchanged by a graded chain homotopy equivalence.  A uniform
homological shift changes its overall sign and a uniform quantum shift
cyclically permutes its coordinates.  Thus neither changes its $\ell^1$
norm.  Taking Euler characteristics after passing to homology proves
\[
 \|\sigma(X)\|_1\leq\dim_{\F} H(X),\qquad
 B_4=\sum_a m_a\,
 \|\sigma(\widetilde{\operatorname{cl}}(C_a))\|_1
 \leq\dim_{\F}\Kh(K;\F).
\]
More precisely, if $E_r,O_r$ are the even and odd homological dimensions
in residue $r$, the deficit is $2\sum_r\min(E_r,O_r)$.
Since reduced Khovanov homology detects the unknot, $B_4>1$ certifies
knottedness.  A value of zero or one is inconclusive.

\paragraph{Aggregation order matters.}
For each whole differential component the implementation recovers relative
quantum shifts using \code{recover\_shifts}.  A nonzero coefficient of dot
degree $d$ between matchings with $m$ arcs and $c$ glued circles requires
$q_{\rm target}-q_{\rm source}=m-c+2d$.  Consistency is checked on all
edges and cycles.  The completed object vectors are shifted, multiplied
by the object's homological sign, and \emph{added before taking a norm}.
Taking separate object norms would discard legitimate differential
cancellations and could produce an unsound rejection.

Exact refinement of a direct-sum decomposition cannot decrease $B_4$:
this is the triangle inequality, together with additivity of the Euler
vector.  Shared isomorphic components add their multiplicities.  The
production scanner caps those multiplicities at three; it preserves
the threshold $\min(2,B_4)$, rather than the full uncapped bound.
Tests compare both exact multiplicities and this saturated threshold.

\subsection{Computing the vector without enumerating the cube}

For a nonempty closed diagram use the raw reduced convention
\[
 J(q)=\sum_{s\in\{0,1\}^n}(-q)^{|s|}
             (q+q^{-1})^{c(s)-1}.
\]
Its quantum degrees have one parity
$\epsilon=(c(0)-1)\bmod2$.  Put $a=J(1)$ and
$d=i^{-\epsilon}J(i)$, which is an integer.  Only residues $\epsilon$
and $\epsilon+2$ can occur, and
\[
 \sigma_\epsilon=(a+d)/2,\qquad
 \sigma_{\epsilon+2}=(a-d)/2.
\]
The engine verifies the parity and the exact divisions.  The existing
orientation-based closure calculation supplies the unreduced value
$2a$ at $q=1$; no exponential state sum is involved.

For the other specialization, construct the actual completed dart
involution and follow the face permutation
$d\mapsto\operatorname{rot}(\alpha(d))$.  Check that its Euler
characteristic is that of a disjoint union of spheres, then checkerboard
color its faces.  A disconnected projection has $J(i)=0$.
For a connected projection choose either color as black, with $v$ black
regions.  At crossing $e$, let $b_e$ be the smoothing bit which does
\emph{not} join its black sectors, put $w_e=(-1)^{b_e}$, and set
$B=\sum_e b_e$.  With $L_G^{(0)}$ a grounded cofactor of the signed
Tait Laplacian,
\[
                  J(i)=(-i)^{B+v-1}\det L_G^{(0)}.
\]
One-circle smoothings correspond to spanning trees of the Tait graph;
expanding their signed weights gives this formula by the matrix-tree
identity.  The identity is polynomial in the edge weights, so negative
weights do not require a positivity assumption.  Tait loops contribute
to $B$ even though they make no Laplacian contribution.  Dropping this
phase, or replacing the determinant by its absolute value, would corrupt
cancellation within a component.

Either checkerboard color is valid if its phase is computed consistently.
The maintained engine chooses the smaller color class to reduce the
dense cofactor dimension.  Fraction-free Bareiss elimination, with row
pivoting and exact integer division, computes the determinant.  Its
integer intermediate sizes are polynomially bounded for these
integer Laplacians; the determinant calculation for one completion has
polynomial bit complexity in its crossing count.

For a single parity vector the norm is
$\max(|a|,|d|)$, using
$|a+d|+|a-d|=2\max(|a|,|d|)$.  Thus the observation retains both Euler
and determinant information.  The significant new possibility is their
application to separate genuine completed components, whose contributions
may cancel in the global polynomial.  Abstract examples of stronger
component bounds do not by themselves establish that a knot scan realizes
those examples.

\subsection{A budgeted observer on the same exact scan}

The driver first tries the existing component Euler obstruction, then
the four-coordinate bound.  Stage zero is always eligible.  Later
observations are skipped while there is just one component copy, since
its completed Euler vector is the unchanged global vector up to shifts.
Observation resumes when a nontrivial direct-sum decomposition appears.
The suffix preparation and completed values are cached by stage and
boundary matching.  Euler and determinant observations share the
\code{euler\_max\_states} allowance, default 4096 distinct queries.

An additional \code{shadow\_max\_work} allowance, default one million,
charges geometry visits, dense matrix allocation and arithmetic updates.
It is an operation allowance, not a count of bit operations or a hard
wall-clock guarantee.  The cooperative global deadline is polled during
geometry and between determinant elimination rows, including on cache
hits.  An incomplete determinant is never cached as a completed value.

If either local inference allowance is exhausted, the driver disables
further observations and continues the \emph{same} capped exact scan.
No restart, discarded completed pivot, or speculative verdict is needed.
Global resource exhaustion still returns \code{UNKNOWN} through the
recognition API.  An unknot verdict is available only when the complete
rank computation has closed.  The empty diagram is handled by its known
unreduced rank two.

This is an adaptive policy in the sense suggested by introsort: expensive
optional work has a budget and a complete fallback.  Its guarantee is
that inference cannot continue consuming unbounded local work after
that allowance.  Unlike introsort's worst-case sorting theorem, it does
not improve the fallback's asymptotic bound.  The default recognition
backend remains unchanged because the measurements below do not show
a broad end-to-end benefit.

\subsection{A reusable terminal kernel, including singular interiors}

Report 26 also proposes factoring a common signed Laplacian into interior
vertices $I$ and boundary terminals.  A completion identifies terminals;
it must leave interior vertices distinct.  A nonsingular interior block
could be eliminated by an ordinary Schur complement, but signed weights
allow a singular interior even when some completed cofactor is nonzero.
The archive uses symmetric one- or two-coordinate pivots to extract a
nonsingular principal block $P$.  The remaining interior block is zero,
of dimension $r$.  With residual coupling $R$, terminal matrix $S$, and
grounded partition matrix $Q$ for a query with $p$ terminal blocks, its
cofactor is
\[
 \det P\;\det
 \begin{pmatrix}
  0&RQ\\ Q^{\mathsf T}R^{\mathsf T}&Q^{\mathsf T}SQ
 \end{pmatrix},\qquad Q\text{ has }q=p-1\text{ columns}.
\]
If $r>q$, or $RQ$ lacks row rank $r$, the answer vanishes.  In
particular, $r>b-1$ makes every query vanish when there are $b$ terminals.
At the critical size $q=r$ the answer becomes
\[
                 (-1)^r\det P\,\det(RQ)^2,
\]
independent of $S$.  Consequently two nonzero critical answers have the
same sign and rational square class; their product is a positive integer
square.  These are identities for a verified common graph and partition,
not a license to identify arbitrary regions of unrelated completions.

\paragraph{An explicit geometric producer.}
The new research program
\path{fast/determinant_research/terminal_geometry_audit.py}
retains the original black regions incident to the suffix crossings,
and retains exactly those crossings as signed edges.  Original black
regions touching a frontier edge become terminals.  For each completed
matching it verifies that completed faces retain the original colors,
that each retained original black region maps to one completed face,
and that every fiber containing more than one old region consists only
of terminals.  These conditions establish the graph quotient by checking
each retained edge's endpoints; no dense matrix comparison is required.
The phase uses the \emph{completed} black vertex count.

The program constructs one archive terminal kernel per suffix and verifies
its algebraic certificate.  Disconnected completed projections have zero
determinant specialization and are handled separately.  Any other failed
geometry check declines the query; it cannot certify a knot.
Across 300 random braid attempts, 78 validated knots and all encountered
prefix matchings, 3186 connected queries passed the checks and agreed
with direct completed determinants.  A further 1317 queries had
disconnected projections.  There were no color, region-splitting, or
interior-identification failures; the largest terminal set had size ten.
The input words, crossing orders, and source hashes accompany these
4503 queries in \code{data/determinant-terminal-geometry-audit.json}.
This is finite evidence for a useful construction, not a proof that every
realizable matching passes its checks.  The recognizer currently uses
direct completed determinants; kernel reuse needs a maintained
implementation, interruptible exact arithmetic, and measurements of
setup amortization before production dispatch is justified.

\subsection{Independent checks and actual performance}

The unchanged delivered package passes its 35 tests.  Its full validation
rerun covers 701 diagrams, 133923 direct cube states, 83 homology
calculations, 62904 columns checked for $d^2=0$, and 4000 quotient queries
on 160 signed graphs, including 57 singular interiors.  The upstream
probe performs 176 scans of 44 diagrams, using two crossing orders and
both physical and shared components.  All 820 proper-prefix observations
obey the exact-rank bound and the appropriate monotonicity condition.

Ten maintained tests independently compare Bareiss elimination with
permutation expansion, completed vectors with direct cube polynomials,
and component bounds with reduced homology.  They check ordinary and
capped multiplicities, mirror diagrams, global deadlines, partial
determinant interruption, shared cache limits, exact fallback with zero
inference budget, invalid inputs, and actual CLI dispatch.  They also
verify that observing a component does not mutate its differential.
The full integrated suite passes all 560 tests.

One concrete strict improvement is the natural crossing order of $T(3,5)$.
Its initial four-coordinate norm is one.  After nine of ten crossings,
the completed components have total norm seven, while their reduced
Euler bound is only one.  The marked backend therefore rejects at stage
nine; the Euler backend reaches the final stage ten.  Similar one-stage
improvements occur on the measured $T(3,7)$ and $T(3,11)$ inputs.
An earlier certified stage does not automatically imply a faster runtime.

\begin{center}\small
\begin{tabular}{@{}lrrrrr@{}}
\toprule
Input & Full/shadow & A/A & Raw/shadow & Raw/Euler & A/A \\
\midrule
unknot & 0.985 & 0.997 & 0.079 & 0.969 & 1.020 \\
trefoil & 1.014 & 1.028 & 3.015 & 0.980 & 0.974 \\
conway & 0.979 & 1.015 & 0.948 & 0.990 & 0.993 \\
kinoshita\_terasaka & 1.007 & 1.013 & 0.951 & 0.992 & 0.987 \\
hard\_unknot\_8 & 0.991 & 1.002 & 0.854 & 1.019 & 1.002 \\
grid\_scrambled\_unknot & 0.986 & 0.982 & 0.803 & 0.980 & 0.979 \\
stress\_braid5\_36 & 1.000 & 0.993 & 360.190 & 1.032 & 0.987 \\
conway\_sum\_2 & 1.004 & 1.004 & 1.475 & 1.539 & 1.007 \\
conway\_sum\_8 & 1.012 & 0.969 & 4.144 & 5.603 & 1.003 \\
torus\_3\_5 & 1.007 & 1.002 & 1.023 & 1.002 & 1.003 \\
torus\_3\_7 & 1.002 & 1.017 & 1.000 & 0.983 & 0.986 \\
torus\_3\_11 & 0.994 & 1.000 & 0.972 & 0.976 & 0.997 \\
\bottomrule
\end{tabular}
\end{center}
Ratios are median paired times; values above one favor the denominator. Full uses the previous default recognizer; raw uses the saturated scanner. A/A compares identical implementations in each scope.


The complete recognizer shows no material benefit from the new backend
on these inputs: old-default/marked ratios range from 0.979 to 1.014,
while the identical-control range is 0.969--1.028.  The changed default
entry point has old/new ratios 0.962--1.019, with median 0.998.
In contrast, the raw marked scanner is about $3.02$ times faster on the
trefoil and $360$ times faster on the 36-crossing stress braid, rejecting
both at stage zero.  Existing recognition filters already decide those
inputs.  The Conway sums gain from component Euler information available
to the old Euler backend too; that backend is faster than the marked
observer on both measured sums.  The raw marked scanner is slower on
Conway, Kinoshita--Terasaka, both nonempty unknot fixtures, and the empty
diagram.  The one-stage torus improvements give ratios of 1.023, 1.000,
and 0.972, respectively, against saturated scanning.  These results
support optional use and further work on selective observation, rather
than enabling the backend by default.
For scale, the empty raw scan rises from about $2.56$ to $34.90$
microseconds, while the stress braid falls from about $635$ to $1.72$
milliseconds; these are medians of the individual arm times.

The reproducer is \code{fast/benchmark\_shadow.py}.  It compares twelve
inputs over seven shuffled paired rounds, with a duplicate control arm.
Recognition measurements keep every default filter enabled and include
fresh validated diagram construction.  A separate raw-decision scope
bypasses those filters and compares the saturated, Euler and marked
scanners.  Each call has a five-second and 50000-object allowance;
any timed repetition returning \code{UNKNOWN} suppresses ratios for
that case.  The historical recognition control loads only
\code{recognize.py} from commit \code{3a6d80ed6}, sharing its unchanged
dependencies.  Imports and one warmup are excluded.  Complete samples
and source hashes are retained in
\code{fast/results/marked\_shadow\_20261008.json}.

\subsection{What is still missing for a general complexity theorem}

Each determinant query is polynomial, but the number and size of physical
relative components need not be small.  The report's stopping-prefix
bound has the conditional form $\operatorname{poly}(n,R)2^{O(W)}$,
where $R$ bounds the actual maintained physical data and $W$ bounds the
frontier through the first decisive observation.  Neither parameter
has a universal bound sufficient for quasi-polynomial recognition.
A planar separator bound on a diagram alone does not bound the surviving
algebraic multiplicities or ensure an early decisive decomposition.

The new observer is useful precisely because it may stop before computing
all homology.  That possibility does not establish that it stops early
on every nontrivial knot, nor that the complete fallback is subexponential
on all unknots.  The earlier assessment of Lackenby's hierarchy claim
remains unchanged: it is a plausible geometric route, while this
implementation still lacks a general quasi-polynomial or subexponential
running-time proof.
